\MakeAbstract{
电子对抗网络中，节点位置会影响探测覆盖、干扰压制和边界任务点性能。为了提升网络的探测和干扰对抗性能，并应对复杂区域中凹边界、空洞及非连通子区域带来的可行位置搜索难题，本文在多目标粒子群优化-区域分解与坐标变换框架上提出了空地协同的节点部署优化方法。

方法采用赫特尔-梅尔霍恩凸分解技术和垂直交点坐标变换法，将粒子编码映射到可行部署区域；通过评估区分空地链路与地地链路路径损耗，以网络覆盖率和最小干扰强度为指标优化节点位置。

仿真结果发现，本文所提算法通过优化节点部署，在网络覆盖和干扰强度指标上较最初配置有明显改善，并在部分场景和部分指标上优于传统算法。本文算法对复杂区域下面向覆盖与干扰压制的节点部署具备实际参考价值。
}{多目标粒子群优化，区域分解，坐标变换，空地协同，电子对抗}

\MakeAbstractEng{
In electronic-countermeasure networks, node locations affect sensing coverage, jamming suppression, and boundary task-point performance. To improve sensing and countermeasure performance while addressing feasible-position search difficulties caused by concave boundaries, holes, and disconnected subregions in complex areas, this thesis proposes an air-ground cooperative node deployment optimization method based on multi-objective particle swarm optimization with decomposition and transformation (MOPSO-DT).

The method uses Hertel-Mehlhorn convex decomposition and a vertical-intersection coordinate transformation to map particle encodings to feasible deployment regions. It distinguishes Air-to-Ground (A2G) and Ground-to-Ground (G2G) path loss in evaluation, and optimizes node positions using network coverage rate and minimum jamming intensity as metrics.

Simulation results show that the proposed algorithm improves network coverage and jamming-intensity metrics over the initial configuration through node deployment optimization, and outperforms traditional algorithms in selected scenarios and on selected indicators. The proposed algorithm provides practical reference for node deployment oriented to coverage and jamming suppression in complex regions.
}{Multi-objective Particle Swarm Optimization, Region Decomposition, Coordinate Transformation, Air-ground Coordination, Electronic Countermeasure}
